\documentclass[10pt,oneside]{book}
%%%% Page Info + Commands %%%%%{

%packages
\usepackage{pgfplots}
\usepackage{mathrsfs}
\pgfplotsset{compat=1.18}
\usepackage{amsmath}
\usepackage{amsfonts}
\usepackage{charter}
\usepackage{amsthm,letltxmacro}
\usepackage{amssymb}
\usepackage{fancyhdr}
\usepackage{float}
\usepackage[most]{tcolorbox}
\usepackage{enumerate}
\usepackage{enumitem}
\usepackage[dvipsnames]{xcolor}
\usepackage{changepage}
\usepackage{graphicx}

% SMALL PAGE COMMAND
\newcommand{\smallpage}{  \setlength \oddsidemargin{.25in}
            \setlength \textwidth{5.5in}
            \setlength \topmargin{0in}
            \setlength \textheight{9in}}


% Commands for sets of numbers
\newcommand{\Z}{\mathbb{Z}}\newcommand{\R}{\mathbb{R}}\newcommand{\N}{\mathbb{N}}\newcommand{\Q}{\mathbb{Q}}\newcommand{\C}{\mathbb{C}}
\newcommand{\real}[1]{\text{Re}\left(#1\right)}
\newcommand{\im}[1]{\text{Im}\left(#1\right)}

% gordie spacing and boxing commands
\newcommand{\gspc}{\vspace{0.13cm}} % 0.5 space
\newcommand{\gline}[0]{\gspc\\} % 1.5 space
\newcommand{\gbline}{\gspc\gspc\\} % double space
\newcommand{\gbox}[1]{\fbox{\parbox{\linewidth}{#1}}} % multiline box command
\newcommand{\gmod}[1]{\,(\text{mod}\,#1)}


\newcommand{\gimage}[2]{
\begin{figure}[H]
    \centering
    \includegraphics[width=#2]{#1}
\end{figure}
}

\newcommand{\centerit}[1]{{\begin{center} #1 \end{center}}}
\newcommand{\ul}[1]{\underline{#1}}

\newcommand{\prt}[2]{\frac{\partial#1}{\partial#2}}

%%%%%%%% command for graphics %%%%%%%%%%%%%

%}

% COLOR BOX

\tcbset{problembox/.style={
    colback=gray!4,
    colframe=gray!50, 
	  boxrule=0pt,
    leftrule=2.5pt, 
    left=2mm, right=2mm, top=1mm, bottom=1mm,
    boxsep=2mm,
    breakable,
    sharp corners
  }
}

\tcbset{subproblembox/.style={
    colback=gray!12,
    colframe=gray!80, 
	  boxrule=0.5pt,
    leftrule=2.5pt, 
    left=2mm, right=2mm, top=1mm, bottom=1mm,
    boxsep=2mm,
    breakable,
    sharp corners
  }
}

% PROBLEM
\newenvironment{problem}[1]
  {\vspace{-0.1cm}\begin{tcolorbox}[problembox]
   \textit{Problem. }#1}
  {\end{tcolorbox}\gspc}

  \newenvironment{subproblem}[2]
  {\vspace{-0.1cm}\begin{tcolorbox}[subproblembox]
   \textit{#1 }#2}
  {\end{tcolorbox}\gspc}


%%%% Page Setup %%%%%%%
\smallpage\sloppy
\pagestyle{fancy}
\lhead{\large{\textbf{PS2 - Maxwell's Equations}}}
\setlength{\headheight}{10pt}


% color commands
\newcommand{\cg}[1]{\color{OliveGreen}#1\color{white}}
\newcommand{\cb}[1]{\color{BlueViolet}#1\color{white}}


\newcommand{\cpr}[1]{\left(#1\right)}
%%%%%%%%%%%%% PHYSICS SPECIFIC COMMANDS

\newcommand{\vA}{\vec{\mathbf{A}}}
\newcommand{\vB}{\vec{\mathbf{B}}}
\newcommand{\vC}{\vec{\mathbf{C}}}
\newcommand{\vZ}{\vec{\mathbf{0}}}
\newcommand{\norm}{\vec{\mathbf{n}}}
\newcommand{\pvec}[1]{\left\langle#1\right\rangle}
\newcommand{\ar}{\vec{\mathbf{r}}}
\newcommand{\arhat}{\hat{\mathbf{r}}}
\newcommand{\vv}{\vec{\mathbf{v}}}
\newcommand{\delfull}{\pvec{\prt{}{x}, \prt{}{y},\prt{}{z}}}
\newcommand{\delinput}[3]{\pvec{\prt{}{x}\cpr{#1}, \prt{}{y}\cpr{#2},\prt{}{z}\cpr{#3}}}
\newcommand{\crossprod}[6]{\text{Det}\left|\begin{matrix}\hat{\textbf{i}}&\hat{\textbf{j}} &\hat{\textbf{k}}\\#1 & #2 & #3\\#4 & #5&#6\end{matrix}\right|}
\newcommand{\dps}{\displaystyle}
\newcommand{\delmatrix}[3]{\begin{bmatrix}\dps\hat x&\dps\hat y&\dps\hat z\\\dps\prt{}x&\dps\prt{}y&\dps\prt{}z\\\dps#1&\dps#2&\dps#3\end{bmatrix}}

%%%%%%%%%%%%%%%%%%%%%%%%%%%%
%%%%%%%%%%%%%%%%%%%%%%%%%%%%%%
%%%%%%%%%%%%%%%%%%%%%%%%%%%%%
%%%%%%%%%%%%%%%%%%%%%%%%%%%%%
%%%%%%%%%%%%%%%%%%%%%%%%%%%%%%
%%%%%%%%% begin doc


\begin{document}


\newcommand{\vE}{{\mathbf{E}}}
\newcommand{\vr}{{\mathbf{r}}}
\newcommand{\etahat}{\hat{\mathbf{\eta}}}
\textbf{Chapter 2 problems:} 6, 7, 12, 15, 18, 22, 29
\subsection*{Problem 6}
Find the electric field a distance $z$ above the center of a flat circular dist of radius $R$ that carries a uniform surface charge $\sigma$. What does your formula give in the limit $R\to\infty$? Also check the case $z\gg R$. 
\begin{problem}
  Via the textbook, we have that the formula for an electric field of surface charge is given by:
  \begin{align*}
    \vE(\vr) = \frac{1}{4\pi \varepsilon_0}\int \frac{\sigma(\vr')}{\eta^2}\hat\eta \,da'
  \end{align*}
  Thus, in order to solve this equation, we note that our integral must cover from $[0,2\pi]$ and from $[0, R]$. Furthermore, we know that in this case, $da' = \sigma d\ell'$.\gline{}
  When we integrate out from 0 to $R$, the charge stored in each ring is equal to $\sigma 2\pi \cdot r$, where $r$ is our integration variable. Thus, in Cylindrical coordinates, we have:
  \begin{align*}
    \frac{1}{4\pi\varepsilon_0}\int_0^{2\pi}\int_0^R \frac{\sigma\cdot 2\pi\cdot \pvec{r\cos\phi, r\sin\phi, -z}}{(r^2 + z^2)^{3/2}}r\, dr d\phi\\
  \end{align*}
  Solving this integral with Mathematica gives us that:
  \begin{align*}
    \vE = \frac{\sigma}{2\varepsilon_0}\cpr{1-\frac{z}{\sqrt{R^2 + z^2}}}
  \end{align*}
  If we observe, as $R\to\infty$, our equation evaluates to $\frac{\sigma}{2\varepsilon_0}$. However, as $z\to\infty$, we get that the electirc field goes to 0. 
\end{problem}
\subsection*{Problem 7}
Find the electric field a distance $z$ from the center of a spherical surface of radius $R$ that carries a uniform charge density $\sigma$. Treat the case $z < R$ (inside) as well as $z>R$ (outside). Express your answers in terms of the total charge $q$ on the sphere.
\begin{problem}
  This problem is somewhat similar to the previous, in that we just take a surface integral to determine the electric field. In this case, we have that:
  \begin{align*}
    r' &= \pvec{z\cos\theta\sin\phi, z\sin\theta\sin\phi,z\cos\phi}\\
    &=\pvec{r, \theta, \phi}\Rightarrow \text{ in spherical.}
  \end{align*}
  However, we're going to simplify this problem by just points the sphere in the direction of the $z$, because this is symmetrical. Thus we get that $z$ is just going to rest at $\pvec{z, 0, 0}$.\gline{} 
  Now, to get the distance to $z$, it's simply just the $\sqrt{R^2 + z^2 + R\cos\phi}$ in each of the spherical directions. Now, with that $dq = \sigma R^2\sin\phi d\phi d\theta$, we get the disgusting integral 
  \begin{align*}
    &=\frac{1}{4\pi\varepsilon_0}\int_0^{2\pi}\int_0^{\pi}\frac{\sigma R^2 (z-R\cos\phi)}{(R^2+z^2-2Rz\cos\phi)^{3/2}}d\phi d\theta\\
    &=\frac{1}{4\pi\varepsilon_0}(2\pi R^2\sigma)\int_0^\pi \frac{(z-R\cos\phi)\sin\phi}{(R^2+z^2-2Rz\cos\theta)^{3/2}}\\
    &\text{Solving using mathematica gives: }\\
    &=\frac{1}{4\pi\varepsilon_0}\frac{2\pi R^2\sigma}{z^2}\cpr{\frac{(z-R)}{|z-R|}-\frac{(z+R)}{|z+R|}}
  \end{align*}
  Via this formulation, if $z > R$, then $\vE = \frac{1}{4\pi\varepsilon_0}\frac{4\sigma\pi R^2}{z^2}$. However, if inside the sphere, because the electric field from all directions will cancel out, we have that $\vE = 0$. 
\end{problem}
\subsection*{Chapter Problem 12}
Use Gauss' law to find the electric field inside and outside a spherical shell of radius $R$ that carries a uniform surface charge density $\sigma$. Compare your answer to Prob. 2.7; notive how much quicker and easier Gauss' method is. 
\begin{problem}
  Ok, this problem will be way simpler because of the niceties of Gauss' law. We know that the surface area of the sphere is $4\pi^2 r^2$. Thus, we get to go like:
  \begin{align*}
    |E|(4\pi r^2) = \frac{Q_\text{enclosed}}{\varepsilon_0}
  \end{align*}
  For $Q_\text{enclosed}$ by the entire surface, we have that it's just $\sigma 4\pi R^2$. Thus, we have that:
  \begin{align*}
    |E| = \frac{Q_\text{enclosed}}{\varepsilon_0 4\pi^2 r^2}=\frac{4\pi R^2\sigma}{4\pi r^2\varepsilon_0} = \frac{\sigma R^2}{r^2\varepsilon_0}\hat r
  \end{align*}
\end{problem}
\subsection*{Problem 15}
Find the electric field inside a sphere that carries a charge density proportional to the distance from the center, $\rho = kr$ for some constant $k$. [\textit{Caution:} This charge density is \textit{not} uniform, and you must \textit{integrate} to get the enclosed charge.]
\begin{problem}
  Ok, for this one, we will first calculate the total charge enclosed for the constant given by $\rho$. 
  \begin{align*}
    Q&=\int_0^\pi\int_0^{2\pi}\int_0^R kr\,r^2\sin\phi\,drd\theta d\phi\\
    &=\int_0^\pi\int_0^{2\pi}\left[\frac{r^4}{4}k\sin\phi\right]\,d\theta d\phi\\
    &= R^4\pi k
  \end{align*}
  Now we apply Gauss' formula:
  \begin{align*}
    |E|A &= \frac{Q_\text{enclosed}}{\varepsilon_0}\Rightarrow A = 4\pi R^2\\
    |E| &= \frac{1}{4\varepsilon_0}kr^2 \hat r
  \end{align*}
\end{problem}
\subsection*{Problem 18}
An infinite plane slab, of thickness $2d$ carries a uniform volume charge density $\rho$. Find the electric field, as a function of $y = 0$ at the center. Plot $E$ versus $y$, calling $E$ positive when $\vE$ point in the $+y$ direction and negative when it points in the $-y$ direction.
\begin{problem}
  \gimage{image.png}{8cm}
\end{problem}
\subsection*{Problem 22}
Find he potential inside and outside a uniformly charge solid sphere whose radius is $R$ and whose total charge is $q$. Use infinity as your reference point. Compute the gradient of $V$ in each region, and check that it yields the correct field. Sketch $V(r)$. 

\begin{problem}
  First note that starting from infinite, we have that:
  \begin{align*}
    V(r)=\int_r^\infty \vE \,d\ell
  \end{align*}
  Very far away from the sphere, we treat it like a point charge, such that:
  \begin{align*}
    \frac{1}{4\pi\varepsilon_0}\frac{q}{r^2}
  \end{align*}
  However, inside the sphere we need Gauss' law:
  \begin{align*}
    \frac{Q}{\frac{4}3\pi R^3}
  \end{align*}
  Substituting in for our electric field gives:
  \begin{align*}
    Q_\text{enclosed} = \cpr{\frac{q}{\frac{4}3\pi R^3}}\cpr{\frac{4}3\pi r^3}\Rightarrow q\frac{r^3}{R^3}
  \end{align*}
  Then, solving for $E$ in Gauss' equationjust gives us:
  \begin{align*}
    \vE = \frac{1}{4\pi\varepsilon_0}\frac{qr}{R^3}\hat r
  \end{align*}
  Sketching this, we get:
  \gimage{image2.png}{8cm}
\end{problem}
\subsection*{Probelm 29}
Use equation 2.29 to calculate the potential inside a uniformly charged solid sphere of radius $R$ and total charge $q$. 
\begin{problem}
  Equation 2.29 is given as:
  \begin{align*}
    V(\vr)=\frac{1}{4\pi\varepsilon_0}\int \frac{\rho(\vr')}{\eta}\,d\tau'
  \end{align*}
  I actually am really lost on how to do this. 
\end{problem}
\end{document}